% authority: science-draft
% status: proposal-only
% route: ontology-law-research-packet
% trigger_classification: derivation_critical_missing_source_law
% adoption_status: blocked_adoption_open_continuation
\documentclass[11pt]{article}
\usepackage[margin=1in]{geometry}
\usepackage{amsmath,amssymb,amsthm}
\usepackage{booktabs}
\usepackage{enumitem}

\newtheorem{definition}{Definition}
\newtheorem{theorem}{Theorem}
\newtheorem{proposition}{Proposition}
\newtheorem{remark}{Remark}
\newcommand{\Ftwo}{\mathbb F_2}
\newcommand{\Disc}{\mathsf{Disc}_{\mathrm{src}}}
\newcommand{\AdmVar}{\mathsf{AdmVar}_{\mathrm{src}}}
\newcommand{\NatCert}{\mathsf{NatCert}_{\mathrm{src}}}
\newcommand{\RobCert}{\mathsf{RobCert}_{\mathrm{src}}}
\newcommand{\EqCand}{\mathsf{EqSrc}^{\mathrm{cand}}}
\newcommand{\Bottom}{\mathsf{Bottom}_{\mathrm{src}}}

\title{Proposal-Only EqSrc Intrinsic Discriminator and Admissibility Law Candidate v1:\\
Chain-Homology Construction, Naturality, Independence, and Robustness}
\author{Ontology Formalizer control packet RT-20260718-027}
\date{2026-07-19}

\begin{document}
\maketitle

\section{Control Status}

This packet defines
\(\mathsf{EqSrcIntrinsicDiscriminatorAdmissibilityLaw}^{\mathrm{cand},v1}_{\mathrm{src}}\)
as a proposal-only source-extension object. Current ontology does not derive
this law, select finite chain complexes as its source packages, or establish
that boundary moves are the physically admissible source variations. Its
status is \texttt{proposal-only},
\texttt{source-extension data}, and
\texttt{blocked\_adoption\_open\_continuation}.

\begin{verbatim}
candidate_result: candidate_formalized_pending_fresh_smuggling_audit
candidate_name: EqSrcIntrinsicDiscriminatorAdmissibilityLaw_src^cand,v1
formal_model: finite_2_truncated_chain_complex_over_F2
intrinsic_discriminator: chi_E:Z_1(E)->H_1(E;F2)
candidate_eqsrc: kernel_pair_of_chi_E
admissible_variations: source_local_boundary_moves_defined_before_chi_E
variation_independence: certified_formally_not_physically
naturality: chain_isomorphism_natural
uniqueness: unique_up_to_intrinsic_codomain_gauge
finite_variation_robustness: proved_for_boundary_move_class
finite_witness: one_vertex_three_loops_one_two_cell
witness_class_count: 4
witness_class_size: 2
current_ontology_derives_candidate: false
physical_admissibility_established: false
candidate_status: proposal-only
adoption_status: blocked_adoption_open_continuation
general_EqSrc_discharged: false
distance_to_gr_ledger_changed: false
freeze_decision: not_frozen
next_route: fresh_smuggling_auditor_review
\end{verbatim}

The exact
\(\mathsf{EqSrcSourceRootedQuotientSelectorCandidate}_{v1}\) construction,
audit, and stress cycle remains locally frozen. The present chain-homology
object is structurally distinct: neither its state space nor its relation is
specified by the prior profile table.

\section{Formal Source Objects, Domains, and Maps}

\begin{definition}[Admitted chain package]
An admitted source package is a finite two-truncated chain complex of
\(\Ftwo\)-vector spaces
\[
 E=\left(C_2(E)\xrightarrow{\partial_2^E}
 C_1(E)\xrightarrow{\partial_1^E}C_0(E)\right),
 \qquad \partial_1^E\partial_2^E=0.
\]
Define the source-state domain, boundary-move subgroup, and intrinsic
codomain by
\[
 A_E=Z_1(E)=\ker\partial_1^E,\qquad
 B_1(E)=\operatorname{im}\partial_2^E,\qquad
 K_E=H_1(E;\Ftwo)=Z_1(E)/B_1(E).
\]
All objects in this definition are source-side finite linear data. No target
atlas, target topology, target metric, detector outcome, benchmark result,
registry, role, validator, handoff, checkpoint, or adoption status occurs in
the definition tree.
\end{definition}

\begin{definition}[Intrinsic discriminator and candidate relation]
For an admitted package \(E\), define
\[
 \Disc(E)=\chi_E:Z_1(E)\longrightarrow H_1(E;\Ftwo),
 \qquad \chi_E(z)=[z].
\]
The proposal-only candidate relation is its kernel pair:
\[
 z\,\EqCand_E\,z'
 \quad\Longleftrightarrow\quad
 \chi_E(z)=\chi_E(z')
 \quad\Longleftrightarrow\quad
 z-z'\in B_1(E).
\]
The codomain is intrinsic to the displayed source complex; it is not a target
classification set or a lookup table imported from desired benchmark
behavior.
\end{definition}

\begin{definition}[Independently generated admissible variations]
Before using \(\chi_E\), its kernel pair, or any selected relation, declare
the elementary source-local boundary moves
\[
 T_b:Z_1(E)\longrightarrow Z_1(E),\qquad
 T_b(z)=z+\partial_2^E b,\qquad b\in C_2(E).
\]
Let \(\AdmVar(E)\) be the submonoid of
\(\operatorname{End}(Z_1(E))\) generated by the maps \(T_b\), with the empty
word as identity and word concatenation as composition. The chain law makes
each \(T_b\) well typed because
\(\partial_1^E(T_bz)=\partial_1^Ez+
\partial_1^E\partial_2^Eb=0\).
\end{definition}

\section{Variation-Class Independence Certificate}

The definition of \(\AdmVar(E)\) uses only
\(C_2(E)\), \(C_1(E)\), \(\partial_2^E\), addition, and finite composition.
It is fixed before \(\chi_E\) and does not reference:
\begin{enumerate}[leftmargin=*]
  \item \(\chi_E\), \(H_1(E)\), \(\EqCand_E\), a kernel pair, or the
  answer-preserving stabilizer of any relation;
  \item a target atlas, target metric, detector protocol, desired exact-GR
  output, or reconstructed target object; or
  \item a role, registry, validation result, handoff, checkpoint, approval, or
  adoption state.
\end{enumerate}
Thus the packet supplies a formal non-circularity certificate:
\[
 \operatorname{Deps}(\AdmVar)
 \cap\{\chi,\EqCand,\operatorname{Stab}(\EqCand),
 \text{target data},\text{workflow authority}\}=\varnothing.
\]
This certificate establishes independence inside the proposed mathematical
model only. It does not establish that the current ontology or nature selects
boundary moves as the physically admissible variation class.

\section{Source-Automorphism Naturality Theorem}

\begin{definition}[Admitted transition]
An admitted transition \(F:E\to E'\) is a chain map
\(F_\bullet=(F_2,F_1,F_0)\) satisfying
\[
 F_1\partial_2^E=\partial_2^{E'}F_2,\qquad
 F_0\partial_1^E=\partial_1^{E'}F_1.
\]
It is a source automorphism when \(E'=E\) and every \(F_i\) is invertible.
A chain isomorphism is an admitted transition with every \(F_i\) invertible.
\end{definition}

\begin{theorem}[Chain-isomorphism naturality]
Every admitted chain map \(F:E\to E'\) restricts to
\(F_1:Z_1(E)\to Z_1(E')\), maps \(B_1(E)\) into \(B_1(E')\), and induces
\[
 H_1(F):H_1(E;\Ftwo)\longrightarrow H_1(E';\Ftwo),
 \qquad H_1(F)[z]=[F_1z].
\]
The discriminator square commutes:
\[
 \boxed{\quad
 \chi_{E'}\circ F_1=H_1(F)\circ\chi_E
 \quad}.
\]
For a chain isomorphism, \(H_1(F)\) is an isomorphism. Consequently every
source automorphism preserves \(\EqCand_E\).
\end{theorem}

\begin{proof}
If \(z\in Z_1(E)\), then
\(\partial_1^{E'}F_1z=F_0\partial_1^Ez=0\), so \(F_1z\) is a cycle. If
\(z=\partial_2^Eb\), then
\(F_1z=\partial_2^{E'}F_2b\), so boundaries map to boundaries. Therefore
\([z]\mapsto[F_1z]\) is well defined. For each \(z\in Z_1(E)\),
\[
 (\chi_{E'}F_1)(z)=[F_1z]
 =(H_1(F)\chi_E)(z),
\]
which proves naturality. An inverse chain map induces the inverse map on
homology. Equality of discriminator values is therefore preserved and
reflected by a chain isomorphism.
\end{proof}

\begin{proposition}[Transition composition and inverse certificates]
For admitted chain maps \(E\xrightarrow{F}E'\xrightarrow{G}E''\),
\[
 H_1(G\circ F)=H_1(G)\circ H_1(F),\qquad H_1(\mathrm{id}_E)=\mathrm{id}_{K_E}.
\]
For a chain isomorphism \(F\),
\(H_1(F^{-1})=H_1(F)^{-1}\). These equalities supply the transition,
composition, and inverse portions of \(\NatCert\).
\end{proposition}

\begin{proof}
Apply both sides to a class \([z]\). Both composites give
\([G_1F_1z]\); identity and inverse follow in the same way.
\end{proof}

\section{Uniqueness up to Intrinsic Codomain Gauge}

\begin{theorem}[Kernel-pair quotient uniqueness]
Let \(D\) be a set and let \(d:Z_1(E)\to D\) be a surjection such that
\[
 d(z)=d(z')\quad\Longleftrightarrow\quad z-z'\in B_1(E).
\]
Then there is a unique bijection
\(\bar d:H_1(E;\Ftwo)\to D\) satisfying
\[
 d=\bar d\circ\chi_E.
\]
Hence any surjective discriminator with exactly the boundary-orbit kernel
pair differs from \(\chi_E\) only by a unique relabeling of its intrinsic
codomain. This relabeling is the permitted source-side codomain gauge.
\end{theorem}

\begin{proof}
Define \(\bar d([z])=d(z)\). If \([z]=[z']\), then
\(z-z'\in B_1(E)\), so \(d(z)=d(z')\); therefore \(\bar d\) is well
defined. Surjectivity follows from that of \(d\). If
\(\bar d([z])=\bar d([z'])\), the kernel-pair hypothesis gives
\([z]=[z']\), so \(\bar d\) is injective. The factorization is immediate.
Any factorization through \(\chi_E\) must send \([z]\) to \(d(z)\), which
proves uniqueness.
\end{proof}

This theorem is conditional on the stated kernel-pair hypothesis. It does not
show that current ontology selects this kernel pair or that every possible
intrinsic discriminator is gauge-equivalent to it.

\section{Finite-Variation Robustness and Sharpness}

\begin{theorem}[Finite boundary-move robustness]
For every finite word
\(V=T_{b_n}\circ\cdots\circ T_{b_1}\in\AdmVar(E)\),
\[
 V(z)=z+\partial_2^E(b_1+\cdots+b_n)
 \quad\text{and}\quad
 \chi_E(V(z))=\chi_E(z)
\]
for all \(z\in Z_1(E)\). Therefore \(\EqCand_E\) is invariant under every
independently admitted finite variation.
\end{theorem}

\begin{proof}
Induction on word length gives the displayed sum, using linearity over
\(\Ftwo\). Its difference from \(z\) lies in \(B_1(E)\), so the quotient
class is unchanged. The empty word is the identity case.
\end{proof}

The certificate \(\RobCert(E)\) consists of the typed generator proof, the
finite-word induction, and the quotient equality above. The robustness claim
is relative to \(\AdmVar(E)\); it is not a claim of robustness under arbitrary
source perturbations.

\section{Concrete Finite Source-Object Witness}

Let \(E_\star\) be the cellular chain package over \(\Ftwo\) with one
0-cell \(v\), three loop 1-cells \(e_0,e_1,e_2\), and one 2-cell \(f\):
\[
 C_0=\Ftwo\langle v\rangle,\quad
 C_1=\Ftwo\langle e_0,e_1,e_2\rangle,\quad
 C_2=\Ftwo\langle f\rangle,
\]
\[
 \partial_1=0,\qquad \partial_2f=e_0+e_1.
\]
Then
\[
 Z_1(E_\star)=\Ftwo^3,\qquad
 B_1(E_\star)=\{0,e_0+e_1\},\qquad
 H_1(E_\star;\Ftwo)\cong\Ftwo^2.
\]
The eight source states form four two-element classes:
\[
\begin{array}{c|c}
\text{class representative}&\text{class}\\\hline
0&\{0,e_0+e_1\}\\
e_0&\{e_0,e_1\}\\
e_2&\{e_2,e_0+e_1+e_2\}\\
e_0+e_2&\{e_0+e_2,e_1+e_2\}.
\end{array}
\]
The nonidentity boundary move \(T_f\) toggles \(e_0+e_1\) and preserves
every class. In contrast, the source translation \(S(z)=z+e_2\) is well
typed on \(Z_1(E_\star)\) but is not generated by boundary moves and changes
every homology class. This is a sharpness witness: the candidate is robust
under its independently declared class but fragile under a concrete
out-of-class variation.

This cell-complex witness is distinct from the frozen four-object quotient
table: it has eight chain states, a boundary subgroup generated by a
2-cell, and a quotient derived from a displayed chain differential.

\section{Fail-Closed Branches}

The proposed law returns \(\Bottom\) and supplies no candidate relation when
any of the following occurs:
\begin{enumerate}[leftmargin=*]
  \item \(\partial_1\partial_2\neq0\), so boundary moves need not preserve
  the cycle domain;
  \item an alleged source state is not in \(Z_1(E)\);
  \item a required nondegenerate witness has trivial or otherwise
  nonseparating \(H_1(E)\);
  \item an alleged transition is untyped, has mismatched chain domains, or
  fails either chain-map square;
  \item transition composition or inverse data are missing or incompatible;
  \item the naturality square fails;
  \item an admitted generator changes \(\chi_E\);
  \item the variation class references \(\chi_E\), \(\EqCand_E\), its
  stabilizer, target data, desired benchmark behavior, or workflow authority;
  \item the quotient factorization or claimed codomain-gauge uniqueness
  certificate fails; or
  \item any target or process datum enters a source-definition dependency
  closure.
\end{enumerate}
Failure closes this candidate instance only. It does not prove a general
EqSrc no-go or future source-extension impossibility.

\section{No-Target and No-Process Guard}

The transitive dependency closure of every carrier, differential, cycle,
boundary, homology class, variation generator, transition, and certificate in
this artifact is confined to displayed finite source algebra. Holding those
source leaves fixed while varying target geometry, benchmark results, or
workflow state leaves every definition and theorem unchanged.

Desired exact-GR recovery is an obligation on any later adopted theory, not a
premise for \(\chi_E\), \(\AdmVar(E)\), or the finite witness. Validator PASS,
registration, and checkpoint status can establish transaction integrity only.

\section{Distance-to-GR, Freeze, and Next Route}

The result is
\texttt{candidate\_formalized\_pending\_fresh\_smuggling\_audit}. It supplies
new definitions, theorems, and a finite witness but does not discharge general
EqSrc and does not change the Distance-to-GR ledger.

The distinct law route is \texttt{not\_frozen}: unlike the locally frozen
profile-table cycle, it provides a new chain-complex object, a naturality
theorem, an independence certificate, a quotient uniqueness theorem, a finite
robustness theorem, and a concrete sharpness witness.

The next lawful packet is one fresh bounded
\texttt{smuggling-auditor@0.2.0} audit of this exact candidate. It must test
whether choosing chain complexes, homology, and boundary moves hides an
unlicensed source primitive; whether the formal independence claim is
overread as physical admissibility; whether target or workflow data enter
transitively; and whether any conclusion extends beyond the stated finite and
conditional scope. It may not repair or adopt the candidate.

\section{Forbidden Conclusions}

This artifact does not authorize canonical ontology modification, source-law
adoption, physical-admissibility establishment, general EqSrc discharge,
RetainH or GenH adoption, \(M_{\mathrm{src}}\) adoption, \(g_{\mathrm{eff}}\)
scope expansion, matter coupling, Einstein equations, finite-variation
robustness beyond the declared class, benchmark promotion, Gate Chair verdict,
completed derivation, future source-extension impossibility, program-level
no-go, or global theory rejection.

\section{Source Materials}

\begin{itemize}[leftmargin=*]
  \item Handoff 0751 and the RT-20260718-026 selector packet.
  \item The registered Ontology Formalizer role contract v0.2.0.
  \item The source-extension classification checklist and no-target-import
  guard map.
  \item Standard finite-dimensional chain-complex, homology, and quotient
  universal-property arguments reproduced explicitly above.
\end{itemize}

\end{document}
